% GENERATED FILE. Edit canonical lessons, then run npm run generate:books.
% canonical-source-hash: fnv1a64:f3cb1f39d706253c
% canonical-lessons: FA-C316-dozd, FA-C316-dozdi, FA-C316-sareq, FA-C316-jib-bor, FA-C316-kif-qap

\chapter{Thieves}
\label{ch:fa-thieves}
\hlchaptermodality{316}

\begin{chapteropening}
\textbf{By the end of this chapter:} \emph{I can talk about thieves and theft: a thief, a theft, a robber, a pickpocket and a bag snatcher.}

Complete the last lesson of chapter 316: I can talk about thieves and theft: a thief, a theft, a robber, a pickpocket and a bag snatcher.
\end{chapteropening}

\section[\texorpdfstring{\fa{دزد} (dozd)}{دزد (dozd)}]{\fa{دزد} (dozd) — a thief}

\label{lesson:FA-C316-dozd}



\begin{quote}
Before the new one: say the Persian for a demonstration, then the Persian for a slogan.
\end{quote}

\subsection*{You'll want to know: \fa{دزد}}
\textbf{\fa{دزد}} (\emph{dozd}) — \textquotedblleft{}a thief\textquotedblright{}.

\begin{quote}
\fa{دزد} \fa{از} \fa{پنجره} \fa{وارد} \fa{خانه} \fa{شد}.
\end{quote}

\begin{quote}
\emph{Dozd az panjere vâred-e khâne shod.}
\end{quote}

\begin{quote}
The thief came into the house through the window.
\end{quote}

Plural \emph{dozd-hâ}. The verb \emph{dozdidan} \textquotedblleft{}to steal\textquotedblright{} is built on it.

\begin{cousinweb}
\emph{dozd} is a native Persian word, usually traced to Middle Persian \emph{duzd}, \textquotedblleft{}thief\textquotedblright{}; EtymDB lists it without a parent.
\end{cousinweb}

\begin{grammarlens}[title={What you've built}]
One more word for telling, explaining and describing at B1.
\end{grammarlens}

\subsection*{Your turn}
\begin{itemize}
\raggedright
  \item \emph{Say it:} \emph{dozd}
  \item \emph{Say it:} \emph{dozd}, once more
  \item \emph{Say it:} a sentence of your own with \emph{dozd}
  \item \emph{Recall:} say the Persian for a demonstration, then the Persian for a slogan, then say \emph{dozd} again
\end{itemize}

\begin{tcolorbox}[breakable,colback=teal!4,colframe=teal!35!black,title={Before you move on}]
What does \fa{دزد} mean? (\textquotedblleft{}A thief\textquotedblright{}.) Say it once more.
\end{tcolorbox}
\section[\texorpdfstring{\fa{دزدی} (dozdi)}{دزدی (dozdi)}]{\fa{دزدی} (dozdi) — a theft, stealing}

\label{lesson:FA-C316-dozdi}



\begin{quote}
Before the new one: say the Persian for a slogan, then the Persian for a thief.
\end{quote}

\subsection*{You'll want to know: \fa{دزدی}}
\textbf{\fa{دزدی}} (\emph{dozdi}) — \textquotedblleft{}a theft, stealing\textquotedblright{}.

\begin{quote}
\fa{دزدی} \fa{ساعت} \fa{دو} \fa{شب} \fa{اتفاق} \fa{افتاد}.
\end{quote}

\begin{quote}
\emph{Dozdi sâ'at-e do-ye shab ettefâq oftâd.}
\end{quote}

\begin{quote}
The theft happened at two in the morning.
\end{quote}

\emph{dozd} \textquotedblleft{}thief\textquotedblright{} + \emph{-i}. In speech \emph{dozdi kardan} is \textquotedblleft{}to steal\textquotedblright{}.

\begin{cousinweb}
\emph{dozdi} adds the native noun suffix \emph{-i} to \emph{dozd} \textquotedblleft{}thief\textquotedblright{}, a word usually traced to Middle Persian \emph{duzd}.
\end{cousinweb}

\begin{grammarlens}[title={What you've built}]
One more word for telling, explaining and describing at B1.
\end{grammarlens}

\subsection*{Your turn}
\begin{itemize}
\raggedright
  \item \emph{Say it:} \emph{dozdi}
  \item \emph{Say it:} \emph{dozdi}, once more
  \item \emph{Say it:} a sentence of your own with \emph{dozdi}
  \item \emph{Recall:} say the Persian for a slogan, then the Persian for a thief, then say \emph{dozdi} again
\end{itemize}

\begin{tcolorbox}[breakable,colback=teal!4,colframe=teal!35!black,title={Before you move on}]
What does \fa{دزدی} mean? (\textquotedblleft{}A theft, stealing\textquotedblright{}.) Say it once more.
\end{tcolorbox}
\section[\texorpdfstring{\fa{سارق} (sâreq)}{سارق (sâreq)}]{\fa{سارق} (sâreq) — a robber, a burglar (formal)}

\label{lesson:FA-C316-sareq}



\begin{quote}
Before the new one: say the Persian for a thief, then the Persian for a theft.
\end{quote}

\subsection*{You'll want to know: \fa{سارق}}
\textbf{\fa{سارق}} (\emph{sâreq}) — \textquotedblleft{}a robber, a burglar (formal)\textquotedblright{}.

\begin{quote}
\fa{پلیس} \fa{به} \fa{دنبال} \fa{سارق} \fa{است}.
\end{quote}

\begin{quote}
\emph{Polis be donbâl-e sâreq ast.}
\end{quote}

\begin{quote}
The police are looking for the robber.
\end{quote}

Formal, in the news; in speech \emph{dozd}. Plural \emph{sâreqân}.

\begin{cousinweb}
\emph{sâreq} is borrowed from Arabic \emph{sāriq} (EtymDB lists the Arabic word), the active participle \textquotedblleft{}one who steals\textquotedblright{}.
\end{cousinweb}

\begin{grammarlens}[title={What you've built}]
One more word for telling, explaining and describing at B1.
\end{grammarlens}

\subsection*{Your turn}
\begin{itemize}
\raggedright
  \item \emph{Say it:} \emph{sâreq}
  \item \emph{Say it:} \emph{sâreq}, once more
  \item \emph{Say it:} a sentence of your own with \emph{sâreq}
  \item \emph{Recall:} say the Persian for a thief, then the Persian for a theft, then say \emph{sâreq} again
\end{itemize}

\begin{tcolorbox}[breakable,colback=teal!4,colframe=teal!35!black,title={Before you move on}]
What does \fa{سارق} mean? (\textquotedblleft{}A robber, a burglar (formal)\textquotedblright{}.) Say it once more.
\end{tcolorbox}
\section[\texorpdfstring{\fa{جیب} \fa{بر} (jib bor)}{جیب بر (jib bor)}]{\fa{جیب} \fa{بر} (jib bor) — a pickpocket}

\label{lesson:FA-C316-jib-bor}



\begin{quote}
Before the new one: say the Persian for a theft, then the Persian for a robber.
\end{quote}

\subsection*{You'll want to know: \fa{جیب} \fa{بر}}
\textbf{\fa{جیب} \fa{بر}} (\emph{jib bor}) — \textquotedblleft{}a pickpocket\textquotedblright{}.

\begin{quote}
\fa{در} \fa{اتوبوس} \fa{یک} \fa{جیب} \fa{بر} \fa{گوشی} \fa{مرا} \fa{برد}.
\end{quote}

\begin{quote}
\emph{Dar otobus yek jib bor gushi-ye marâ bord.}
\end{quote}

\begin{quote}
On the bus a pickpocket took my phone.
\end{quote}

\emph{jib} \textquotedblleft{}pocket\textquotedblright{} + \emph{bor}, the present stem of \emph{boridan} \textquotedblleft{}to cut\textquotedblright{}: a pocket-cutter.

\begin{cousinweb}
\emph{jib} \textquotedblleft{}pocket\textquotedblright{} is borrowed from Arabic \emph{jayb} (EtymDB). \emph{bor-} is the present stem of the native verb \emph{boridan} \textquotedblleft{}to cut\textquotedblright{}.
\end{cousinweb}

\begin{grammarlens}[title={What you've built}]
One more word for telling, explaining and describing at B1.
\end{grammarlens}

\subsection*{Your turn}
\begin{itemize}
\raggedright
  \item \emph{Say it:} \emph{jib bor}
  \item \emph{Say it:} \emph{jib bor}, once more
  \item \emph{Say it:} a sentence of your own with \emph{jib bor}
  \item \emph{Recall:} say the Persian for a theft, then the Persian for a robber, then say \emph{jib bor} again
\end{itemize}

\begin{tcolorbox}[breakable,colback=teal!4,colframe=teal!35!black,title={Before you move on}]
What does \fa{جیب} \fa{بر} mean? (\textquotedblleft{}A pickpocket\textquotedblright{}.) Say it once more.
\end{tcolorbox}
\section[\texorpdfstring{\fa{کیف} \fa{قاپ} (kif qâp)}{کیف قاپ (kif qâp)}]{\fa{کیف} \fa{قاپ} (kif qâp) — a bag snatcher}

\label{lesson:FA-C316-kif-qap}



\begin{quote}
Before the new one: say the Persian for a robber, then the Persian for a pickpocket.
\end{quote}

\subsection*{You'll want to know: \fa{کیف} \fa{قاپ}}
\textbf{\fa{کیف} \fa{قاپ}} (\emph{kif qâp}) — \textquotedblleft{}a bag snatcher\textquotedblright{}.

\begin{quote}
\fa{کیف} \fa{قاپ} \fa{کیف} \fa{زن} \fa{را} \fa{گرفت} \fa{و} \fa{دوید}.
\end{quote}

\begin{quote}
\emph{Kif qâp kif-e zan râ gereft va david.}
\end{quote}

\begin{quote}
The bag snatcher grabbed the woman's bag and ran.
\end{quote}

\emph{kif} \textquotedblleft{}bag\textquotedblright{} + \emph{qâp}, from \emph{qâpidan} \textquotedblleft{}to snatch\textquotedblright{}.

\begin{cousinweb}
\emph{kif qâp} joins \emph{kif} \textquotedblleft{}bag\textquotedblright{}, whose origin is not settled, and \emph{qâp-}, the stem of \emph{qâpidan} \textquotedblleft{}to snatch\textquotedblright{}, which is usually traced to Turkic \emph{qap-} \textquotedblleft{}to grab\textquotedblright{}. EtymDB lists no entry.
\end{cousinweb}

\begin{grammarlens}[title={What you've built}]
One more word for telling, explaining and describing at B1.
\end{grammarlens}

\subsection*{Your turn}
\begin{itemize}
\raggedright
  \item \emph{Say it:} \emph{kif qâp}
  \item \emph{Say it:} \emph{kif qâp}, once more
  \item \emph{Say it:} a sentence of your own with \emph{kif qâp}
  \item \emph{Recall:} say the Persian for a robber, then the Persian for a pickpocket, then say \emph{kif qâp} again
\end{itemize}

\begin{tcolorbox}[breakable,colback=teal!4,colframe=teal!35!black,title={Before you move on}]
What does \fa{کیف} \fa{قاپ} mean? (\textquotedblleft{}A bag snatcher\textquotedblright{}.) Say it once more.
\end{tcolorbox}
